\section{Symmetry: the translate pigeonhole}\label{sec:symmetry}

The flag bound cannot hold rotations with interval masks to the pairing rate (\cref{sec:order-says}). This section
does, by a different argument: a code larger than half the space meets its own translate, and on the intersection a
shared decoder forces the data to be translation-covariant.

\subsection{The lemma}

For operators $S,T$ write $[S,T]=ST-TS$.

\begin{lemma}[Pigeonhole against a translate]\label{lem:pigeonhole}
\formal{Pigeonhole.lean}{two_mul_finrank_ker_sub_le_ker,finrank_ker_commutator_sub_le,two_mul_finrank_ker_sub_le,two_mul_finrank_ker_sub_le_of_commute,finrank_range_commutator_le}
Let $\sw,K,T$ be linear operators on a space of finite dimension $N$ over a field. Then
\[
  \dim\ker(\sw-K)\ \le\ \tfrac12\bigl(N+\dim\ker[\sw-K,T]\bigr)\ \le\ \tfrac12\bigl(N+\dim\ker[\sw,T]+\rank[K,T]\bigr).
\]
\end{lemma}

\begin{proof}
Let $\X=\ker(\sw-K)$, $D=\dim\X$ and $\X_T=\{X\in\X:TX\in\X\}$. The map $\X\to F^N/\X$, $X\mapsto TX+\X$, has kernel
$\X_T$ and a target of dimension $N-D$, so $\dim\X_T\ge2D-N$; $T$ need not be invertible. For $X\in\X_T$,
$(\sw-K)TX=0=T(\sw-K)X$, so $\X_T\subseteq\ker[\sw-K,T]$. That is the first inequality. For the second, $[\sw-K,T]X=0$
means $[\sw,T]X=[K,T]X\in\im[K,T]$, and the preimage of a subspace $W$ under a map has dimension at most that of its
kernel plus $\dim W$.
\end{proof}

With $S=\sw-K$ the first inequality says $\rank[S,T]\le2\rank S$: a commutator has rank at most twice the rank of
either factor. The Lean proof takes this route, from $\rank[S,T]\le\rank(ST)+\rank(TS)$ and rank--nullity
(\lean{finrank_range_commutator_le}). The lemma is elementary, and we claim no novelty for it. What is used below is the form with $\sw$ and $K$ separated: $K$
may nearly commute with a $T$ with which $\sw$ does not.

\begin{proposition}[For the swap, the best case is the pairing value]\label{prop:swap-floor}
\partformal{Swap.lean}{card_le_finrank_ker_swap_commutator,two_mul_finrank_ker_swap_sub_one,two_mul_finrank_range_swap_sub_one}{the
equality case for $T=1\otimes G$.}
On $V\otimes V$ with $\dim V=d$, every operator $T$ has $\dim\ker[\sw,T]\ge d$, in every characteristic. Equality holds
for $T=1\otimes G$, that is $X\mapsto XG^{\tr}$, when $G$ is non-derogatory.
\end{proposition}

\begin{proof}
The kernel of $\sw-1$ is the space of symmetric matrices, of dimension $d(d+1)/2$ in every characteristic, so
$\rank(\sw-1)=d(d-1)/2$, and $[\sw,T]=[\sw-1,T]$ has rank at most $d(d-1)$ by \cref{lem:pigeonhole}. For $T=1\otimes G$,
$[\sw,T]X=GX^{\tr}-X^{\tr}G^{\tr}$, so the kernel is the space of $Y=X^{\tr}$ with $GY=YG^{\tr}$. As $G^{\tr}$ is
similar to $G$, this space is isomorphic to the centralizer of $G$, whose dimension is $d$ exactly when $G$ is
non-derogatory.
\end{proof}

So \cref{lem:pigeonhole} never gives less than $(d^2+d)/2$ for the swap, and it gives exactly the pairing value when
$K$ commutes with a suitable~$T$. The rest of the section is about $K$ that commute with a translation up to a
boundary.

\subsection{Class-uniform rules}

Let $\sigma$ be the cyclic rotation of $F^{\Z/d}$, $(\sigma x)_j=x_{j+1}$. For a matrix $X$ with rows
$x_c\in F^{\Z/d}$, $c\in\Z/d$, write $\sigma^\delta X$ for the matrix with rows $\sigma^\delta x_c$.

\begin{definition}[Class-uniform rules]\label{def:class-uniform}
\formal{ClassUniform.lean}{rot,IsClassUniformCode}
Fix a partition of the row index $\Z/d$ into non-empty classes. A \emph{class-uniform rule} is a family of linear maps
$\cA_J\colon\Mat_d(F)\to F^{\Z/d}$, one for each class~$J$. A subspace $\X\subseteq\Mat_d(F)$ is a \emph{code} of the
rule if
\[
  x_c=\cA_J(\sigma^cX)\qquad\text{for every }X\in\X,\text{ every class }J\text{ and every }c\in J .
\]
\end{definition}

In words: row $c$ of $X$ is a fixed linear function of the whole matrix with its columns rotated by~$c$, and the function
depends only on the class of~$c$. The receiver of row $c$ looks at the stored word from its own position, through a
window of any size, and receivers in one class use the same decoder.

In the notation of \cref{sec:onehop} (this translation is not formalized): let $K=\sum_aA_a\otimes B_a$ have second
factors in the span of the matrices $P_J\sigma^e$, a rotation followed by the mask $P_J$ of a class, and arbitrary first
factors, say $K=\sum_{J,e}A_{J,e}\otimes P_J\sigma^e$. Then $(KX)_{r,c}=\sum_{i,e}(A_{J,e})_{r,i}X_{i,c+e}$ for $c\in J$,
so the swap code of $K$ is the largest code of the rule $\cA_J(Y)_r=\sum_{i,e}(A_{J,e})_{r,i}Y_{i,e}$, and every rule
arises in this way. This is the family of \cref{sec:order-says}. For a single class the $B_a$ are circulant and
\cref{cor:commutative} applies. If the mask is applied before the rotation, $K^{\tr}$ has second factors
$P_J\sigma^{-e}$, and since $\rank(\sw-K)=\rank(\sw-K^{\tr})$ the bounds below hold for that family as well.

\begin{theorem}[Theorem~TP, the translate pigeonhole]\label{thm:TP}
\formal{ClassUniform.lean}{IsClassUniformCode.finrank_le,Link,IsClassUniformCode.two_mul_finrank_le,IsClassUniformCode.two_mul_finrank_le_of_const}
Let $\X$ be a code of a class-uniform rule. Let $\Delta_0\subseteq\Z/d$ be a set of $s+1\ge2$ shifts, and let $\chi$
be the number of connected components of the graph on $\Z/d$ in which $c\sim c'$ when $c$ and $c'$ lie in one class and
$c'-c\in\Delta_0-\Delta_0$. Then
\[
  \dim\X\ \le\ \frac{s\,d^2+d\,\chi}{s+1}.
\]
In particular, if the classes are $M$ cyclic intervals, then $\dim\X\le(d^2+Md)/2$.
\end{theorem}

\begin{outline}
Let $\X'=\{X:\sigma^\delta X\in\X\text{ for all }\delta\in\Delta_0\}$, an intersection of $s+1$ preimages of~$\X$;
each has codimension at most $d^2-\dim\X$, so $\dim\X'\ge(s+1)\dim\X-s\,d^2$
(\lean{card_mul_finrank_add_le_iInf_comap}, \lean{mul_finrank_le_add_finrank_iInf_comap}, in
\leanfile{Pigeonhole.lean}). Let $X\in\X'$, $\delta,\delta'\in\Delta_0$ and $c,c'$ in one class $J$
with $c+\delta=c'+\delta'$. The rule applied to $\sigma^\delta X$ at row $c$ gives
$\sigma^\delta x_c=\cA_J(\sigma^{c+\delta}X)$, and applied to $\sigma^{\delta'}X$ at row $c'$ it gives
$\sigma^{\delta'}x_{c'}=\cA_J(\sigma^{c'+\delta'}X)$. The right sides agree, so $x_{c'}=\sigma^{c'-c}x_c$
(\lean{IsClassUniformCode.row_eq_of_rot_mem}). Hence each row of $X$ determines the rows in its component, restriction
to one row per component is injective on $\X'$, and $\dim\X'\le d\chi$. The Lean statement takes any set $S$ of rows
that meets every component and gives $(s+1)\dim\X\le s\,d^2+d\card S$. For $M$ cyclic intervals and $\Delta_0=\{0,1\}$ the components are the intervals. In Lean this case
(\lean{IsClassUniformCode.two_mul_finrank_le}) takes a non-empty set $S$ containing every row whose class differs from
that of the row before it, so that $\card S=M$ for $M\ge2$ intervals and $S$ is any single row for one class
(\lean{IsClassUniformCode.two_mul_finrank_le_of_const}).
\end{outline}

The space on which the proof lands is worth naming. For one unit shift it is
$\Hk=\{X:x_{c+1}=\sigma x_c\text{ whenever }c,c+1\text{ lie in one class}\}$: inside each interval the matrix is constant
along anti-diagonals. It has $d$ free parameters per interval, whatever the length of the interval.

We say that a partition into $M$ classes has an \emph{almost-period} if some single shift $\delta$ leaves only $M$
components in the graph of \cref{thm:TP} with $\Delta_0=\{0,\delta\}$, so that the bound is $(d^2+Md)/2$. Intervals have
the almost-period~$1$, and residue classes modulo $M$ have the almost-period~$M$.

\begin{lemma}[Extremal codes contain $\Hk$]\label{lem:extremal}
\notformalized
Let the classes be $M$ cyclic intervals and let $\X$ be a code of a class-uniform rule with $2\dim\X=d^2+Md$. Then
$\X\cap\sigma^{-1}\X=\Hk$; in particular $\Hk\subseteq\X$.
\end{lemma}

\begin{informalproof}
By the proof of \cref{thm:TP}, $\X'=\X\cap\sigma^{-1}\X$ lies in $\Hk$ and has dimension at least
$2\dim\X-d^2=Md=\dim\Hk$.
\end{informalproof}

\begin{proposition}[Equal pieces: the bound is attained]\label{prop:block-symmetric}
\formal{BlockSymmetric.lean}{blockSymmetric,isClassUniformCode_blockSymmetric,two_mul_finrank_blockSymmetric,blockSymmetric_attains_bound}
Let $d=Mm$ with $M,m\ge1$ and let the classes be the $M$ intervals of length~$m$. The \emph{block-symmetric} matrices,
those whose $m\times m$ blocks in the $M\times M$ grid are all symmetric, are a code of a class-uniform rule, of
dimension $(d^2+Md)/2$.
\end{proposition}

\begin{outline}
Write a row index as $c=sm+k$ and a column index as $r=tm+l$ with $0\le k,l<m$. Block symmetry says
$X_{sm+k,\,tm+l}=X_{sm+l,\,tm+k}$. The right side is the entry of $\sigma^cX$ in row $sm+l$ and column
$tm+k-c=(t-s)m$, which depends on $c$ only through its block~$s$. So $x_c=\cA_s(\sigma^cX)$ with
$\cA_s(Y)_{tm+l}=Y_{sm+l,\,(t-s)m}$ (\lean{isClassUniformCode_blockSymmetric}). There are $M^2$ blocks, each with
$m(m+1)/2$ free entries, and $M^2m(m+1)/2=(d^2+Md)/2$ (\lean{two_mul_finrank_blockSymmetric}); the set of the $M$
first rows of the blocks satisfies the hypothesis of \cref{thm:TP} for the unit shift
(\lean{blockSymmetric_attains_bound}).
\end{outline}

The operator $K$ of this rule is the coarse transposition of the grid: $X\mapsto(\cA_{s(c)}(\sigma^cX))_c$, read as a
matrix with these columns, has the entry $X_{sm+l,\,tm+k}$ in position $(tm+l,\,sm+k)$, so block $(s,t)$ is moved to
position $(t,s)$ with its inner coordinates kept.

So the largest code of a class-uniform rule with $M$ equal interval pieces has dimension exactly $(d^2+Md)/2$, the
value of the $M\times M$ coarse transposition, against the bound $1-1/(2M)$ that flags give for the same family. For
$M=d$ the pieces are single rows, the rule is unrestricted and the code is everything. In between, each additional
piece raises the bound by $d/2$.

\begin{remark}[Numerical: what the exact value does not say]\label{rem:extremal}
Exhaustive computations over $\F_2$ show the following.
(a) Extremal codes are far from unique: for $d=4$ and classes $\{0,1\},\{2,3\}$ there are $59\,392$ codes of
dimension~$12$, in $327$ orbits under the evident symmetries, and the block-symmetric code is one orbit.
(b) Coding beats copying inside the bound: for $d=4$ and pieces $\{0\},\{1,2,3\}$ the bound $12$ is attained, but no
rule that copies one stored entry to each output entry exceeds~$11$.
(c) The bound is not attained for every partition into unequal intervals: by \cref{lem:extremal} and a solver, the
pieces $\{0\},\{1\},\{2,3,4\}$ at $d=5$ have largest code $19<20$ (\cref{rem:small}).
\end{remark}

\begin{proposition}[Copy rules, every partition]\label{prop:copy}
\notformalized
Call a class-uniform rule a \emph{copy rule} if every output entry is a multiple of one stored entry:
$X_{c,r}=\lambda_{J,r}\,X_{\pi_J(r),\,\rho_J(r)+c}$ for $c\in J$. A code of a copy rule with $M$ classes, of any shape,
has dimension at most $(d^2+Md)/2$.
\end{proposition}

\begin{informalproof}
Let $f(c,r)=(\pi_J(r),\rho_J(r)+c)$ for $c\in J$, a map on the $d^2$ positions. A matrix in the code is determined, on
each component of the functional graph of $f$, by its value at one position on the cycle of that component, so the
dimension is at most the number of cycles. A fixed point $(c,r)$ has $\pi_J(r)=c$ with $J$ the class of $c$, so for each
$r$ there is at most one fixed point per class: at most $Md$ in all. Every other cycle has at least two positions.
\end{informalproof}

So for copy rules the shape of the classes does not matter, while for coded rules \cref{thm:TP} needs the classes to
have an almost-period. The difference is real: \cref{prop:aperiodic} exhibits two classes with no almost-period and a
coded rule whose code has dimension greater than $(d^2+Md)/2$. How far coded rules on such classes can go is
\cref{prob:aperiodic}.

\subsection{On a line}

A tape is a line, not a cycle. Let rows and columns be indexed by $\{0,\dots,d-1\}$ and let the classes be intervals. A
\emph{two-sided rule on a line} is given by coefficients $\cA_{J,i}[r][e]$, $-d<e<d$, with
\[
  X_{c,r}=\sum_{i}\sum_{j=0}^{d-1}\cA_{J,i}[r][j-c]\,X_{i,j}\qquad(c\in J).
\]
The decoder of row $c$ again sees the stored word from position $c$, through a window that is cut off at the ends
instead of wrapping around. Every class-uniform rule is of this form, with the same classes: take
$\cA_{J,i}[r][e]$ to depend on $e$ modulo~$d$ (this remark is not formalized). On a line, an interval is a set of
consecutive rows that does not wrap around.

\begin{theorem}[The line form of Theorem~TP]\label{thm:line}
\formal{Line.lean}{IsLineUniformCode,IsLineUniformCode.two_mul_finrank_le,isLineUniformCode_of_coeff,lineView,truncShift}
A code of a two-sided rule on a line with $M$ interval classes has dimension at most $(d^2+Md)/2$.
\end{theorem}

\begin{outline}
Let $(TX)_{i,j}=X_{i,j+1}$ for $j<d-1$ and $(TX)_{i,d-1}=0$, and $\X'=\{X\in\X:TX\in\X\}$, so that
$\dim\X'\ge2\dim\X-d^2$. Let $X\in\X'$ and let $c,c+1$ lie in one class~$J$. The rule for $TX$ at row $c$ reads
$X_{c,r+1}=\sum_i\sum_{j\ge1}\cA_{J,i}[r][j-(c+1)]X_{i,j}$, with $X_{c,d}$ read as~$0$, and the rule for $X$ at row
$c+1$ has the same sum with $j\ge0$. So
\[
  X_{c+1,r}-X_{c,r+1}=\sum_i\cA_{J,i}[r][-(c+1)]\,X_{i,0}.
\]
Suppose that column $0$ of $X$ vanishes. Then inside each class, with rows $a,\dots,a+m_J-1$, the matrix is constant
along anti-diagonals. An anti-diagonal that meets column~$0$ inside the class is zero, and so is one whose first
position in the class is in the last column below row~$a$, because $X_{c,d}=0$. The remaining anti-diagonals start in
row~$a$ at a column $r\ge m_J$: there are $d-m_J$ of them. So
$\dim\X'\le d+\sum_J(d-m_J)=Md$. In Lean the decoder of a class is an arbitrary linear function of the zero-padded
window $(X_{i,e+c})_{i,e}$ (\lean{lineView}); the coefficient form above is an instance
(\lean{isLineUniformCode_of_coeff}). The classes need not be intervals there: the bound is $(d^2+d\card S)/2$ for any
set $S$ of rows that contains row~$0$ and every row whose class differs from that of the row before it. The key step is
\lean{IsLineUniformCode.row_succ_eq}.
\end{outline}

With one class this gives the pairing value for second factors in the span of the truncated shifts, the Toeplitz
matrices, which do not commute and are not covered by \cref{cor:commutative}; the symmetric matrices attain it
(\lean{isLineUniformCode_symmetric}). For $M$ equal intervals the block-symmetric matrices attain the bound on a line
as well, by the remark above; this is not formalized.

\begin{remark}[Numerical: on a line, other shifts fail]\label{rem:line-fails}
For $d=3$ and classes $\{0,2\},\{1\}$, \cref{thm:TP} with $\Delta_0=\{0,2\}$ gives $7$ on a cycle. On a line the largest
code of a two-sided rule has dimension~$8$ (exhaustive over $\F_2$); one such code is $\{X:X_{0,0}=X_{2,2}\}$. So
\cref{thm:line} does not extend to shift sets other than $\{0,1\}$.
\end{remark}

\subsection{Twists, digits and composition}

A rule need not commute with a translation to be caught; it is enough that it intertwines it with boundedly many others.

\begin{theorem}[Twist classes]\label{thm:twist}
\formal{Twist.lean}{succ_mul_finrank_swapCode_le_of_twist_classes,two_mul_finrank_swapCode_le_of_twist,twistSpace,succ_mul_finrank_swapCode_le,finrank_le_sum_finrank_twistSpace}
Let $G\in\Mat_d(F)$, let the terms of $K=\sum_aA_a\otimes B_a$ fall into classes $1,\dots,s$, and let
$G'_1,\dots,G'_s$ commute pairwise, with $GB_a^{\tr}=B_a^{\tr}G'_\kappa$ for every $a$ in class~$\kappa$. Then
\[
  \dim\ker(\sw-K)\ \le\ \frac{s\,d^2+\sum_\kappa h_\kappa}{s+1},\qquad h_\kappa=\dim\{P:G^{\tr}P=PG'_\kappa\}.
\]
\end{theorem}

\begin{outline}
Let $\X'=\{X:XG^t\in\ker(\sw-K),\ 0\le t\le s\}$, of dimension at least $(s+1)\dim\ker(\sw-K)-s\,d^2$
(\lean{succ_mul_finrank_swapCode_le}). For $X\in\X'$ and $P=X^{\tr}$ the hypothesis gives
$(G^{\tr})^tP=\sum_\kappa T_\kappa G_\kappa'^{\,t}$ for $0\le t\le s$, with $T_\kappa=\sum_{a\in\kappa}A_aXB_a^{\tr}$. Applying
$P\mapsto G^{\tr}P-PG'_\kappa$ removes class~$\kappa$ from such a system and lowers $s$ by one, and its kernel is the
space counted by $h_\kappa$; induction on the number of classes gives $\dim\X'\le\sum_\kappa h_\kappa$
(\lean{finrank_le_sum_finrank_twistSpace}). This is the statement that $\prod_\kappa(L-R_\kappa)P=0$ for the commuting
operators $L$, left multiplication by $G^{\tr}$, and $R_\kappa$, right multiplication by $G'_\kappa$.
\end{outline}

\begin{corollary}[Modulation frequencies]\label{cor:weyl}
\notformalized
Let $F$ contain a primitive $d$-th root of unity $\omega$ and let $Z=\mathrm{diag}(1,\omega,\dots,\omega^{d-1})$. If
every $B_a^{\tr}$ is a product $Z^kC$ with $C$ circulant and $k$ in a fixed set of $s$ frequencies, then
$\dim\ker(\sw-K)\le\frac s{s+1}(d^2+d)$.
\end{corollary}

\begin{informalproof}
$\sigma Z=\omega Z\sigma$, so $\sigma Z^kC=Z^kC\,(\omega^k\sigma)$: take $G=\sigma$ and $G'_k=\omega^k\sigma$, which
commute. The matrices $\sigma^{\tr}=\sigma^{-1}$ and $\omega^k\sigma$ are diagonalizable with the same simple spectrum,
the $d$-th roots of unity, so $h_k=d$.
\end{informalproof}

One frequency is the pairing rate, and all $d$ frequencies span the full matrix algebra. The mask of an interval
contains many frequencies (all of them when its length is prime to~$d$), which is why interval masks needed
\cref{thm:TP}.

\begin{corollary}[Every address digit must be broken]\label{cor:digit}
\partformal{Pigeonhole.lean}{two_mul_finrank_ker_sub_le_of_commute}{everything specific to the statement: that $K$
commutes with~$T$, the dimension count $\dim\ker[\sw,T]=d^2/q$, and the sharpness. Lean has only the abstract
lemma.}
Write the index set as $\Z/m\times\Z/q$, $d=mq$, so that the second coordinate is one digit of the address. Suppose
every $B_a$ has the form $B'_a\otimes C_a$ with all $C_a$ in $F[G]$ for one non-derogatory $q\times q$ matrix $G$, for
example with all $C_a$ circulant. Then
\[
  \dim\ker(\sw-K)\ \le\ \tfrac12\Bigl(d^2+\frac{d^2}q\Bigr),
\]
with equality for the partial transpose that exchanges the first digits of the two indices and leaves the second
digits in place.
\end{corollary}

\begin{informalproof}
$K$ commutes with $T=1\otimes(1_m\otimes G)$, and as in \cref{prop:swap-floor} $\ker[\sw,T]$ is isomorphic to the
centralizer of $1_m\otimes G$, which is $\Mat_m(F)\otimes F[G]$, of dimension $m^2q=d^2/q$. Apply
\cref{lem:pigeonhole}. The partial transpose is $\sw_m\otimes1$ on $(F^m\otimes F^m)\otimes(F^q\otimes F^q)$; its second
factors are $E_{u'u}\otimes1_q$, and its swap code is $\Mat_m(F)\otimes\{\text{symmetric }q\times q\text{ matrices}\}$,
of dimension $m^2q(q+1)/2$.
\end{informalproof}

The transposition is the product of the commuting exchanges of the digits of the two addresses
(\cref{sec:saturation}). \Cref{cor:digit} says that a data path which treats any one digit of one address uniformly
is at rate $\tfrac12+\tfrac1{2q}$, however it treats the others; by \cref{lem:pigeonhole}, rate $1-\eps$ needs
$\rank[K,T]\ge(1-2\eps-\tfrac1q)d^2$ for the rotation $T$ of every digit of either index. This is a requirement at
every scale.

\begin{remark}[Composition, and why this gives no logarithm]\label{rem:compose-TP}
\Cref{thm:TP} is a statement about one hop: the class-uniform rules with $M$ classes are not closed under
composition, and two of them compose to a rule with up to $d$ classes. \Cref{lem:pigeonhole} does compose. Since
$[K_2K_1,T]=K_2[K_1,T]+[K_2,T]K_1$, $\rank[K_2K_1,T]\le\rank[K_1,T]+\rank[K_2,T]$
(\lean{finrank_range_commutator_comp_le}), and by induction
\[
  \rank[K_t\cdots K_1,T]\ \le\ \sum_i\rank[K_i,T].
\]
For $T=1\otimes\sigma$ and a hop whose second factors lie in the span of the $P_J\sigma^e$ with $M$ interval classes,
$[K_i,T]X$ is supported on the $M$ columns at which the class changes, so $\rank[K_i,T]\le Md$ and $t$ such hops give
$\dim\ker(\sw-K_t\cdots K_1)\le(d^2+d+tMd)/2$.

The same subadditivity shows why \cref{cor:digit} does not produce a factor $\log d$. An exact transposer needs
$\sum_i\rank[K_i,T]\ge d^2-d^2/q$ for the rotation $T$ of every digit. But let $U$ be the perfect unshuffle in base~$q$
on one index, which moves every digit to the next position and is a linear-time operation on tapes, and let $G_j$ be
the rotation of digit~$j$ of that index. Then $UG_j=G_{j+1}U$, so $[U,G_j]=(G_{j+1}-G_j)U$ has rank $d-d/q$, and
$\rank[1\otimes U,1\otimes G_j]=d^2-d^2/q$ for every~$j$ at once: one operation supplies the whole requirement for all
digits of that index. The requirement is real at every scale and can be met at all scales together.
\end{remark}

\subsection{What the theorem says}\label{sec:symmetry-says}

\begin{itemize}[leftmargin=1.5em]
\item \emph{For translation-covariant rules, the half in the half bound is a pigeonhole.} A code larger than half the
space meets its translate in at least $2D-d^2$ dimensions, and translation-covariant data has $d$ dimensions per uniform
piece. The argument does not depend on the characteristic, on the degree, on the window of the decoder, or on the other
tensor side. It does depend on exact covariance under a unit shift inside each class, on the classes being the same for
every output row, on linearity, and on there being one hop. The flag bounds of \cref{sec:order} are not shown to be
instances of it.
\item \emph{No symmetry that the swap lacks.} By \cref{lem:pigeonhole}, a code that exceeds the pairing value by
$\eps d^2$ needs $\rank[K,T]+\dim\ker[\sw,T]\ge d+2\eps d^2$ for every operator~$T$. Together with \cref{sec:order}:
to beat the pairing rate, the decoder may have neither a hidden order nor a symmetry, on either index, that the swap
does not share.
\item \emph{It uses the rule, not the relation.} The hypothesis of \cref{thm:TP} is that translated receivers share a
decoder. That is the feature of machines singled out in \cref{sec:rule}, and it is invisible in the cone relation.
\item \emph{The folklore step is optimal in dimension.} With $M$ equal pieces, nothing beats the $M\times M$ coarse
transposition by even one dimension, although many other codes equal it (\cref{rem:extremal}).
\end{itemize}
